\section{Resumen del Caso de Estudio}

El caso se centra en el transporte diario (domicilio-trabajo o viceversa) de profesionales de la salud en turnos AM y PM en Santiago de Chile. El área de logística de la empresa de transporte debe diseñar rutas óptimas considerando los nodos origen--destino de cada profesional y sus horarios de turno.



%La organización busca planificar y optimizar el transporte diario de sus profesionales hacia y desde sus lugares de trabajo. Para lograrlo, el área de logística debe diseñar diariamente las rutas para distintos turnos, considerando las ubicaciones geográficas, las ventanas de tiempo requeridas y el estricto cumplimiento de normativas sindicales de comodidad.



%\section{Objetivos y complejidad del problema}
%El objetivo principal es desarrollar un modelo que genere rutas óptimas para minimizar los costos operacionales y la distancia recorrida. La complejidad radica en satisfacer restricciones sindicales de cuatro categorías de pasajeros, donde solo se permite transportar simultáneamente a personas de categorías consecutivas, además de cumplir con tiempos máximos de espera y viaje.







%Más allá del desafío puramente operacional de enrutar una flota y minimizar los costos de transporte , este problema reviste una profunda relevancia práctica para la estrategia de la empresa. El diseño de un sistema de transporte altamente confiable y puntual impacta de manera directa en el bienestar y la calidad de vida de los profesionales de la salud. Al disminuir la carga mental asociada al traslado diario y reducir significativamente las tasas de ausencias o retrasos, la empresa asegura la continuidad operativa de sus centros médicos, lo que se traduce de forma inmediata en una mejor calidad del servicio y atención entregada a los pacientes.

%Adicionalmente, resolver este problema logístico genera impactos corporativos tangibles. Por ejemplo, permite a la organización fortalecer su reputación institucional, retener talento al ofrecer mejores condiciones laborales y consolidar su posición en el mercado de la salud. Por último, la optimización eficiente de los vehículos compartidos alinea a la empresa con metas de sostenibilidad y responsabilidad social, reduciendo la huella de carbono de sus operaciones y mitigando su impacto en la congestión vehicular urbana.










%\section{Objetivos y Complejidad del Problema}

%El objetivo es desarrollar un modelo que genere rutas óptimas minimizando costos operacionales, definidos por tipo/cantidad de vehículos a utilizar y su distancia recorrida. El problema exige respetar estrictas restricciones temporales y sindicales, incluyendo tiempos máximos de espera y compatibilidad entre categorías de pasajeros.



\section{Relevancia Práctica e Impacto Real}

La resolución de este problema logístico trasciende la optimización matemática de rutas. Representa una decisión estratégica que fortalece directamente la propuesta de valor y competitividad de la empresa a través de tres dimensiones fundamentales:

%POSIBLE RESUMEN DE LO ANTERIOR POR TEMAS DE ESPACIO
%la \textbf{continuidad operativa}, asegurando que el personal crítico llegue puntualmente a sus turnos; el \textbf{bienestar laboral}, al transformar el transporte en un beneficio que reduce el estrés del personal médico; y la \textbf{eficiencia corporativa}, optimizando el uso de recursos y disminuyendo costos operacionales.



\begin{itemize}
    \item \textbf{Continuidad operativa y calidad asistencial:} En el sector de la salud, la puntualidad del personal es un factor crítico. Un sistema de transporte confiable mitiga directamente los riesgos asociados a ausencias o retrasos. Al garantizar que los profesionales lleguen a sus turnos de manera oportuna, la empresa asegura la continuidad de sus operaciones y, en consecuencia, mantiene un alto estándar en la calidad del servicio entregado a los pacientes.
    
    \item \textbf{Bienestar laboral y retención de talento:} El cumplimiento riguroso de las normativas sindicales respecto a los tiempos de viaje no es solo una restricción del modelo, sino una herramienta de cuidado al trabajador. Optimizar el traslado disminuye significativamente la carga mental y el estrés de los profesionales de la salud. Esto se traduce en un mejor clima laboral, fortaleciendo la reputación institucional y posicionando a la empresa como un lugar atractivo para retener talento clave.
    
    \item \textbf{Sostenibilidad y eficiencia corporativa:} Desde una perspectiva organizacional, consolidar los viajes compartidos de manera eficiente permite cuidar el presupuesto al reducir costos operacionales innecesarios. Paralelamente, esta eficiencia alinea a la empresa con sus objetivos de responsabilidad social y ambiental, ya que una menor cantidad de vehículos en tránsito contribuye activamente a la reducción de la huella de carbono y alivia la congestión del tráfico urbano en Santiago.
\end{itemize}


\section{Propuesta}
Para materializar este valor, se propone el desarrollo de un \textbf{Sistema de Soporte a la Toma de Decisiones (DSS)} programado en Python. Esta herramienta operará como un \textit{pipeline integral} que procesará automáticamente las solicitudes de viaje, incorporando un modelo predictivo de tráfico y una metaheurística de ruteo (ALNS) para generar hojas de ruta óptimas en una hora o menos. Con esto, se ofrece una solución técnica que automatiza y profesionaliza la planificación diaria de la flota.



%Se propone una metodología que permita ejecutar el modelo para que los viajes compartidos sean competitivos y eficientes, optimizando rutas y la asignación de clientes, cumpliendo con las restricciones y objetivos planteados, clasificando el caso como un problema relevante de optimización de rutas de transporte de personal.

%Para materializar aquello, el entregable final consistirá en un Sistema de Soporte a la Toma de Decisiones (DSS) desarrollado en Python. Este leerá y procesará las solicitudes de viajes e integrará el modelo predictivo de tráfico con una metaheurística de enrutamiento para procesar las restricciones en un tiempo límite de una hora, generando automáticamente hojas de ruta detalladas para los despachadores de la flota. De esta forma, se ofrece una solución integral que automatiza y optimiza la planificación diaria del transporte.